\documentclass[11pt,a4paper]{article}

\usepackage[romanian]{babel}
\usepackage[utf8]{inputenc}
\usepackage[T1]{fontenc}
\usepackage{lmodern}
\usepackage{amsmath,amssymb,amsthm}
\usepackage{enumitem}
\usepackage{booktabs}
\usepackage{geometry}
\geometry{margin=2.5cm}

\title{Niveluri de izolare ale tranzacțiilor în baze de date}
\author{}
\date{}

\begin{document}
\maketitle

\section{Introducere}

Într-un sistem de gestiune a bazelor de date (SGBD) multi-utilizator, mai multe tranzacții rulează concurent și accesează aceleași date. Pentru a păstra consistența, SGBD-ul trebuie să controleze modul în care aceste tranzacții se „văd” una pe alta. Acest control se face prin \emph{nivelurile de izolare} ale tranzacțiilor.

Izolarea este una dintre proprietățile ACID:
\begin{itemize}[leftmargin=1.5cm]
  \item \textbf{Atomicitate} -- tranzacția este indivizibilă: ori se execută complet, ori nu se execută deloc.
  \item \textbf{Consistență} -- tranzacția duce baza de date dintr-o stare consistentă în alta.
  \item \textbf{Izolare} -- efectele unei tranzacții nu sunt vizibile altor tranzacții „în mijlocul” execuției.
  \item \textbf{Durabilitate} -- odată confirmate, modificările supraviețuiesc căderilor de sistem.
\end{itemize}

Standardul SQL definește patru niveluri de izolare:
\begin{enumerate}[leftmargin=1.5cm]
  \item Read Uncommitted
  \item Read Committed
  \item Repeatable Read
  \item Serializable
\end{enumerate}

Unele SGBD-uri introduc și variante bazate pe versiuni (de exemplu, \emph{Snapshot Isolation}), dar acestea se raportează tot la aceleași fenomene de concurență.

\section{Fenomenelor de concurență}

Nivelurile de izolare sunt definite în funcție de ce tipuri de anomalii sunt permise sau interzise.

\subsection{Dirty read}

\textbf{Dirty read} apare când o tranzacție citește date modificate de o altă tranzacție care nu a făcut încă \texttt{COMMIT}.

Exemplu:
\begin{itemize}[leftmargin=1.5cm]
  \item T1: \texttt{UPDATE Cont SET sold = sold - 100 WHERE id = 1;} (fără \texttt{COMMIT} încă)
  \item T2: citește \texttt{sold} pentru contul 1 și vede valoarea diminuată.
  \item T1: face \texttt{ROLLBACK}.
\end{itemize}
T2 a citit o valoare care, logic, „nu a existat niciodată” în baza de date confirmată.

\subsection{Non-repeatable read}

\textbf{Non-repeatable read} apare când o tranzacție citește de două ori același rând și obține valori diferite, deoarece o altă tranzacție a modificat acel rând între cele două citiri.

Exemplu:
\begin{itemize}[leftmargin=1.5cm]
  \item T1: citește \texttt{sold} pentru contul 1.
  \item T2: modifică \texttt{sold} pentru contul 1 și face \texttt{COMMIT}.
  \item T1: citește din nou \texttt{sold} pentru contul 1 și vede o altă valoare.
\end{itemize}

\subsection{Phantom read}

\textbf{Phantom read} apare când o tranzacție execută de două ori aceeași interogare pe un set de rânduri, iar între cele două execuții o altă tranzacție inserează sau șterge rânduri care satisfac aceeași condiție.

Exemplu:
\begin{itemize}[leftmargin=1.5cm]
  \item T1: \texttt{SELECT * FROM Comenzi WHERE total > 1000;}
  \item T2: inserează o comandă nouă cu \texttt{total > 1000} și face \texttt{COMMIT}.
  \item T1: re-execută interogarea și obține un set de rânduri diferit (apare un „fantom”).
\end{itemize}

\subsection{Serialization anomaly}

\textbf{Serialization anomaly} apare când rezultatul unui set de tranzacții concurente nu este echivalent cu niciun ordin serial (una câte una) al acelor tranzacții. Este o formă mai generală de anomalie, care poate implica combinații de citiri și scrieri.

\section{Tabel de sinteză pentru nivelurile de izolare}

Standardul SQL leagă nivelurile de izolare de fenomenele permise sau interzise:

\begin{center}
\begin{tabular}{lcccc}
\toprule
\textbf{Nivel} & \textbf{Dirty read} & \textbf{Non-repeatable read} & \textbf{Phantom read} & \textbf{Serialization anomaly} \\
\midrule
Read Uncommitted & permis & permis & permis & permis \\
Read Committed   & interzis & permis & permis & permis \\
Repeatable Read  & interzis & interzis & permis & permis \\
Serializable     & interzis & interzis & interzis & interzis \\
\bottomrule
\end{tabular}
\end{center}

\section{Read Uncommitted}

\subsection{Descriere}

\textbf{Read Uncommitted} este cel mai slab nivel de izolare. Tranzacțiile pot citi date neconfirmate ale altor tranzacții. Acest lucru maximizează concurența, dar permite toate tipurile de anomalii: dirty read, non-repeatable read, phantom read.

\subsection{Avantaje și dezavantaje}

\begin{itemize}[leftmargin=1.5cm]
  \item \textbf{Avantaje:} blocări minime, throughput mare, potrivit pentru rapoarte unde o mică inexactitate este acceptabilă.
  \item \textbf{Dezavantaje:} consistență slabă, risc de decizii bazate pe date care vor fi anulate.
\end{itemize}

\subsection{Exemplu SQL (generic)}

\begin{verbatim}
SET TRANSACTION ISOLATION LEVEL READ UNCOMMITTED;
BEGIN TRANSACTION;

SELECT * FROM Conturi;

COMMIT;
\end{verbatim}

\section{Read Committed}

\subsection{Descriere}

\textbf{Read Committed} este, de obicei, nivelul implicit în multe SGBD-uri. O tranzacție poate citi doar date care au fost deja confirmate de alte tranzacții. Dirty read este eliminat, dar non-repeatable read și phantom read pot apărea.

\subsection{Avantaje și dezavantaje}

\begin{itemize}[leftmargin=1.5cm]
  \item \textbf{Avantaje:} elimină citirile murdare, oferă un echilibru bun între consistență și performanță.
  \item \textbf{Dezavantaje:} aceeași interogare poate da rezultate diferite în cadrul aceleiași tranzacții.
\end{itemize}

\subsection{Exemplu SQL}

\begin{verbatim}
SET TRANSACTION ISOLATION LEVEL READ COMMITTED;
BEGIN TRANSACTION;

SELECT sold FROM Conturi WHERE id = 1;

-- între timp, altă tranzacție poate modifica sold-ul și face COMMIT

SELECT sold FROM Conturi WHERE id = 1;

COMMIT;
\end{verbatim}

\section{Repeatable Read}

\subsection{Descriere}

\textbf{Repeatable Read} garantează că, dacă o tranzacție citește un rând, nici o altă tranzacție nu poate modifica sau șterge acel rând până când prima tranzacție nu se termină. Astfel, non-repeatable read este eliminat. Totuși, pot apărea phantom read, deoarece alte tranzacții pot insera rânduri noi care satisfac aceeași condiție de căutare.

\subsection{Avantaje și dezavantaje}

\begin{itemize}[leftmargin=1.5cm]
  \item \textbf{Avantaje:} citirile asupra acelorași rânduri sunt consistente în cadrul tranzacției.
  \item \textbf{Dezavantaje:} mai multe blocări, risc mai mare de blocaje (deadlock), phantom read încă posibil.
\end{itemize}

\subsection{Exemplu SQL}

\begin{verbatim}
SET TRANSACTION ISOLATION LEVEL REPEATABLE READ;
BEGIN TRANSACTION;

SELECT * FROM Conturi WHERE id = 1;

-- altă tranzacție nu poate modifica/șterge rândul cu id = 1
-- dar poate insera un cont nou care să apară în alte interogări

COMMIT;
\end{verbatim}

\section{Serializable}

\subsection{Descriere}

\textbf{Serializable} este cel mai puternic nivel de izolare. Execuția concurentă a tranzacțiilor este garantată a fi echivalentă cu o execuție serială (una după alta) într-o anumită ordine. Sunt eliminate dirty read, non-repeatable read și phantom read, precum și anomaliile de serializare.

\subsection{Avantaje și dezavantaje}

\begin{itemize}[leftmargin=1.5cm]
  \item \textbf{Avantaje:} consistență maximă, comportament ușor de raționat (ca și cum tranzacțiile ar rula pe rând).
  \item \textbf{Dezavantaje:} blocări extinse, concurență redusă, risc crescut de blocaje și time-out-uri.
\end{itemize}

\subsection{Exemplu SQL}

\begin{verbatim}
SET TRANSACTION ISOLATION LEVEL SERIALIZABLE;
BEGIN TRANSACTION;

SELECT * FROM Comenzi WHERE total > 1000;

-- altă tranzacție nu poate insera/șterge/ modifica rânduri
-- care ar afecta rezultatul interogării de mai sus

COMMIT;
\end{verbatim}

\section{Snapshot Isolation (versiuni de rânduri)}

Unele SGBD-uri (de exemplu, SQL Server, PostgreSQL) oferă un mecanism de izolare bazat pe versiuni de rânduri, numit adesea \textbf{Snapshot Isolation}. Ideea este că fiecare tranzacție vede o „fotografie” (snapshot) consistentă a bazei de date la momentul începerii tranzacției, fără a bloca citirile.

\subsection{Caracteristici}

\begin{itemize}[leftmargin=1.5cm]
  \item Citirile nu blochează scrierile și invers, deoarece se lucrează cu versiuni.
  \item Elimină dirty read și non-repeatable read.
  \item Phantom read poate fi eliminat în funcție de implementare.
  \item Pot apărea conflicte la commit (de exemplu, două tranzacții modifică același rând pe baza aceleiași versiuni).
\end{itemize}

\section{Compromisul între consistență și concurență}

Nivelurile de izolare reprezintă un compromis între:
\begin{itemize}[leftmargin=1.5cm]
  \item \textbf{Consistență} (mai puține anomalii, comportament mai „serial”)
  \item \textbf{Concurență} (mai puține blocări, throughput mai mare)
\end{itemize}

În practică:
\begin{itemize}[leftmargin=1.5cm]
  \item Read Uncommitted este rar folosit în aplicații critice.
  \item Read Committed este un default rezonabil pentru multe aplicații.
  \item Repeatable Read și Serializable se folosesc în zone unde consistența este mai importantă decât performanța.
  \item Snapshot Isolation oferă adesea un compromis bun: citiri consistente cu blocări reduse.
\end{itemize}

\section{Concluzii}

Nivelurile de izolare sunt esențiale pentru proiectarea și înțelegerea comportamentului tranzacțiilor într-o bază de date. Alegerea nivelului potrivit depinde de:
\begin{itemize}[leftmargin=1.5cm]
  \item cerințele de consistență ale aplicației,
  \item toleranța la anomalii de citire,
  \item cerințele de performanță și scalabilitate.
\end{itemize}

Înțelegerea fenomenelor (dirty read, non-repeatable read, phantom read, serialization anomaly) și a modului în care fiecare nivel de izolare le permite sau le interzice este cheia pentru a configura corect un SGBD și pentru a scrie tranzacții robuste.

\end{document}
